\documentclass[11pt]{article}
\usepackage[a4paper, margin=2.5cm]{geometry}
\usepackage{amsmath, amssymb, amsthm, amsfonts}
\usepackage{mathtools}
\usepackage{enumitem}
\usepackage{hyperref}
\usepackage{tcolorbox}
\usepackage{longtable}
\usepackage{booktabs}
\usepackage{array}

\setlength{\emergencystretch}{3em}
\sloppy

\newtheorem{theorem}{Theorem}[section]
\newtheorem{lemma}[theorem]{Lemma}
\newtheorem{proposition}[theorem]{Proposition}
\newtheorem{corollary}[theorem]{Corollary}
\newtheorem{conjecture}[theorem]{Conjecture}
\newtheorem{hypothesis}[theorem]{Hypothesis}
\theoremstyle{definition}
\newtheorem{definition}[theorem]{Definition}
\newtheorem{remark}[theorem]{Remark}
\newtheorem{example}[theorem]{Example}
\newtheorem{interpretation}[theorem]{Interpretive Claim}

\newcommand{\R}{\mathbb{R}}
\newcommand{\N}{\mathbb{N}}
\newcommand{\Q}{\mathbb{Q}}
\newcommand{\C}{\mathbb{C}}
\newcommand{\Z}{\mathbb{Z}}
\newcommand{\OO}{\mathbb{O}}
\newcommand{\HH}{\mathbb{H}}
\newcommand{\pphi}{\varphi}
\newcommand{\Ocal}{\mathcal{O}}
\newcommand{\Vsub}{V_{600}}
\newcommand{\Cphi}{C_{\pphi}}
\newcommand{\Fix}{\operatorname{Fix}}
\newcommand{\Schur}{\operatorname{Schur}}

\title{Cascade Closure: A Conditional Compatibility Template\\
for Geometric State Selection\\[0.4em]
\large From Objective Reduction to Crystallisation\\
under Projection Constraint}
\author{Lee Smart \\ Institute of Vibrational Field Dynamics}
\date{2026-05}

\begin{document}

\maketitle

\begin{abstract}
We propose a single named candidate mechanism, formulated as a formal
\emph{compatibility template} for propagating a supplied
closure-compatible fixed point through a projection stack of
geometric carriers. The mechanism does not by itself select that
fixed point; it specifies the conditions under which residual
compatibility and gradient-flow fixed-point status compose downward in
systems where unresolved geometric mismatch propagates through a
layered closure substrate:
\emph{cascade closure}. Penrose Objective Reduction is the
neighbouring geometry-coupled state-selection model for which this
paper is proposed as a structural analogue, rather than as a
refutation or as a strict generalisation (we do not derive
Objective Reduction as a special case of cascade closure); \emph{collapse}
becomes the external description of the resolution event,
\emph{crystallisation} the internal description, and \emph{cascade}
the propagation process. We formalise the mechanism as a rung-indexed
closure-residual functional with an explicit projection stack, prove
a conditional residual-compositionality proposition, identify the binary
icosahedral 600-cell vertex graph $\Vsub$ as the chosen finite carrier
of the construction (without claiming uniqueness), and document one
architecture-level mapping to a computational system (the ARIA
observer-process tuple
$\Ocal = \langle B, S, M, G, I, C, A\rangle$) and one passive physical
witness (the $\Vsub$ closure-response operator $\Cphi = L_M + \pphi^{-2} I$
with $\|\Cphi^{-1}\| = \pphi^2$ exactly). The mathematical payload of
the paper is a definition plus a conditional propagation lemma; we
do not exhibit a non-trivial instance in which all clauses of the
cascade-closure event occur. Instead, in lieu of a worked occurrence,
we supply a \emph{consumer API} for successor papers --- a list of
the four explicit obligations (O0)--(O3) a successor paper must
discharge to invoke cascade closure as a load-bearing premise ---
and we record the status of those obligations in the most
load-bearing imported tower in the programme (the finite-level P3/P4
refinement spaces plus closure dynamics). With the source-side
companion paper \cite{CascadeRefinementCompat}, the obligations
discharge as follows: (O0) satisfied at the finite-dimensional
level; (O3) closed in an abstract scalar pure-midpoint refinement
model for the scaled curvature operator
$\widetilde A_n = 2^n A_n^{\mathrm{vertex}}$ via an exact
Schur-complement identity; (O2) \emph{as stated} (operator/flow
level) remains false, with the boundary-energy / Schur substitute
discharged in the same abstract model; (O1) reduces to
finite-dimensional spectral decomposition on the connected graph
Laplacian, with nontrivial selection out of scope. Lifting the (O3)
discharge and the (O2) substitute from the abstract model to the
full P3/P4 tower of \cite{CascadeRefinementSpaces, CascadeClosureDynamics}
is conditional on two open hypotheses (L1) and (L3) of
\cite[\S\texttt{sec:scope-and-gap}]{CascadeRefinementCompat}; the third lift hypothesis
(L2) is already a theorem of P3. This makes the paper the programme's
named dependency gate, not its selection-mechanism endpoint. The paper is a mechanism
paper, not a programme paper. It is not a replacement for quantum
mechanics, not a derivation of consciousness, not a cosmogenesis
theorem, and not a proof of any open mainstream-mathematics problem.
\end{abstract}

\begin{tcolorbox}[title=\textbf{Epistemic status (read first)}]
\begin{itemize}[leftmargin=*]
\item The cascade-closure mechanism (Definition~\ref{def:cascade-closure})
  is the paper's organising definition, not a theorem of physics.
\item The compositionality proposition (Proposition~\ref{prop:resid-comp})
  is conditional on two named hypotheses.
\item No nontrivial example is shown here to satisfy all clauses of
  Definition~\ref{def:cascade-closure}; the mathematical payload is
  a definition plus a conditional propagation lemma, not a
  demonstration that the mechanism occurs.
\item In lieu of a worked occurrence, the paper supplies a
  \emph{consumer API} (\S\ref{ssec:consumer-api}): the explicit
  four obligations (O0)--(O3) a successor paper must discharge to
  invoke cascade closure. The status of these obligations in the
  imported P3/P4 finite-level tower is recorded in
  \S\ref{ssec:worked-invocation}; with the source-side companion
  \cite{CascadeRefinementCompat}, (O3) is closed only in an
  abstract scalar pure-midpoint refinement model, (O2) \emph{as
  stated} remains false (a boundary-energy / Schur substitute is
  discharged in the same model), and (O1) reduces to finite-dimensional
  spectral decomposition with nontrivial selection out of scope.
  P3/P4-tower lift of the (O3) discharge and (O2) substitute is
  conditional on the open hypotheses (L1) and (L3) of
  \cite[\S\texttt{sec:scope-and-gap}]{CascadeRefinementCompat}.
\item The binary icosahedral 600-cell vertex graph $\Vsub$ is the chosen
  finite carrier; this paper does not prove uniqueness.
\item ARIA is mapped to the observer-process \emph{architecture} at
  the process-description level; we do not claim ARIA realises the
  full cascade-closure event of Definition~\ref{def:cascade-closure}
  (clauses (1)--(3)), that ARIA is biologically conscious, or that it
  proves qualia.
\item The closure-response operator $\Cphi$ has an external substrate
  witness (the \(b\to s\mu^+\mu^-\) angular anomaly fit), used here
  as one short external-use example, not as a model-selection claim.
\item Standalone mass-spectrum, brain/EEG, visible-universe, and
  Millennium-class applications (beyond the one-paragraph cited
  ARIA and $b$-anomaly substrate witnesses used here) are out of
  scope for this paper. Each belongs in its own paper.
\end{itemize}
\end{tcolorbox}

\begin{tcolorbox}[title=\textbf{Reader map}]
\begin{itemize}[leftmargin=*]
\item \S\ref{sec:motivation}: Penrose OR and the state-selection problem.
\item \S\ref{sec:vocabulary}: collapse, reduction, cascade, crystallisation
  --- precise vocabulary for the mechanism.
\item \S\ref{sec:mechanism}: the formal mechanism: rung-indexed
  closure-residual functional, cascade-closure event, conditional
  residual-compositionality proposition, and the consumer API
  (\S\ref{ssec:consumer-api}--\S\ref{ssec:worked-invocation}) by
  which successor papers invoke cascade closure.
\item \S\ref{sec:carrier}: the chosen finite carrier $\Vsub$ and its
  closure-response operator $\Cphi$.
\item \S\ref{sec:aria}: ARIA as an architecture-level observer-process
  mapping for the tuple.
\item \S\ref{sec:witness}: one externally reported substrate witness
  ($b$-anomaly closure kernel).
\item \S\ref{sec:discussion}: discussion, claim ledger, what remains
  out of scope.
\item Appendix~\ref{app:field-signature-audit}: re-runnable empirical
  anchor for the load-bearing data claims.
\item Appendix~\ref{app:aria-runtime-map}: ARIA runtime module mapping
  to the observer-process tuple.
\item Appendix~\ref{app:rung-placement}: programme placement of
  cascade closure (compressed nine-row navigation table for
  successor-paper authors).
\end{itemize}
\end{tcolorbox}

\section{Motivation: Penrose OR and the state-selection problem}\label{sec:motivation}

The standard quantum-mechanical vocabulary calls the resolution of an
indeterminate state into a definite outcome \emph{wavefunction
collapse} or \emph{measurement}. This vocabulary does not by itself
specify a substrate-level selection mechanism. The Born rule supplies
probabilities; the projection postulate specifies the update rule.

Penrose~\cite{Penrose1996} proposed that gravitational self-energy of
a mass-energy superposition supplies a substrate-coupled threshold
above which state reduction occurs objectively, independent of any
observer. The Penrose--Hameroff extension~\cite{PenroseHameroff1996}
situates this Objective Reduction within biological microtubules. The
core structural move is geometry-coupled state selection: a
substrate-supplied condition --- not an observer --- produces the
reduction. We take Penrose OR as the immediate conceptual neighbour of
the present work. We do not target it for refutation, and we do not
adopt its specific gravitational threshold; we abstract one
structural form.

The abstraction is the following. Replace ``one substrate-coupled
threshold'' with ``a stratified hierarchy of geometric closure layers,
each supplying its own residual functional.'' Replace the
``superposition-exceeds-threshold then reduction'' picture with one
in which an unresolved closure residual at one layer is modelled as
propagating, under the hypotheses below, through admissible
projection maps to a closure-compatible fixed point on a lower layer.
The shape of the question is the same as Penrose's: state selection
coupled to substrate geometry. The mechanism inside is different.

This paper formalises that mechanism. We do not propose a replacement
for quantum mechanics. We do not claim to derive measurement
probabilities; the Born rule remains the input. We do not claim to
solve the measurement problem; we define a candidate mechanism for
state selection in systems satisfying the stated closure and
projection hypotheses.

\section{Vocabulary}\label{sec:vocabulary}

The mechanism rests on a small precise vocabulary.

\begin{description}[leftmargin=2.5em, style=nextline]
\item[Collapse] the external description of loss of alternatives ---
the Born-rule statement that, after the event, only one outcome
remains accessible to the observer. Collapse names the event from
outside.
\item[Reduction] the physical state-selection event itself, considered
without committing to a particular ontology. Penrose Objective
Reduction is the canonical example; the present cascade-closure
construction is proposed as a further substrate-geometry
parametrisation.
\item[Cascade] the propagation of an unresolved closure residual
through the admissible projection maps of a stratified substrate.
The cascade is the dynamics between instability onset and
crystallisation.
\item[Crystallisation] the internal description of the resolution
event: selection of a closure-compatible limit state on the
substrate. Crystallisation names the event from inside.
\end{description}

Within this vocabulary, the four terms name complementary aspects of the same proposed event-type from different
viewpoints: \emph{reduction} is the event; \emph{collapse} is its
external view; \emph{cascade} is its propagation dynamics on a
stratified substrate; \emph{crystallisation} is its internal
fixed-point-selection content.

\section{Formal mechanism}\label{sec:mechanism}

\subsection{The closure-residual functional}\label{ssec:closure-residual}

Let the substrate be a finite (possibly stratified) sequence of
\emph{rungs} indexed by $k \in \{0, 1, \ldots, N\}$. At each rung the
substrate exposes a geometric carrier $G_k$ (a graph, polytope,
algebra, or projection class) equipped with a closure-residual
functional
\begin{equation}\label{eq:Rk}
   R_k : \Phi_k \longrightarrow [0, \infty),
\end{equation}
where $\Phi_k$ is the space of field-state configurations on $G_k$. For a configuration $\phi_k \in \Phi_k$, the functional
measures the failure of $\phi_k$ to satisfy the rung-$k$ closure
condition encoded by $R_k$. We say $\phi_k$ is
\emph{residual-compatible} (at rung $k$) when $R_k(\phi_k) = 0$.
This is the only sense in which we use the word
``closure-compatibility'' as a property of a single state at a single
rung. The fully assembled cascade-event compatibility (zero residual,
flow invariance, terminal projection-compatibility across the whole
tower) is the separate notion stipulated in
Definition~\ref{def:cascade-closure} below.

\paragraph*{Standing analytic hypotheses.} Throughout this section we
assume that, at each rung $k$, the configuration space $\Phi_k$ is a
finite-dimensional smooth manifold \emph{without boundary} equipped
with a Riemannian metric (boundary, constrained, or
infinite-dimensional Hilbert-manifold cases are admissible
extensions, but require additional analytic preconditions stated below
or analogues thereof; see also Remark~\ref{rem:clause2-redundancy}
for one place where ``without boundary'' is load-bearing); that
$R_k \in C^1(\Phi_k)$ so that the gradient $\nabla R_k$ is defined
with respect to the chosen Riemannian metric; that the
negative-gradient flow $\Psi_k^t$ generated by $-\nabla R_k$ exists
and is unique for all $t \geq 0$ from each initial condition needed
below; and that each pushforward $\pi^{\sharp}_{k+1, k}$ is
$C^1$ as a map $\Phi_{k+1} \to \Phi_k$. These are the analytic
preconditions for the gradient-flow language used in
Definition~\ref{def:cascade-closure} below; they are not asserted to
hold automatically and are part of the per-instance verification.

The finite-dimensional closure-flow infrastructure used as background
here --- gradient operator, energy dissipation, conditional coercive
contraction, mass-block compatibility, flow existence, and mass-only
flow intertwining --- is established in \cite{CascadeClosureDynamics}
for the source's quadratic closure functional
$F_n = \alpha R_n + \beta E_n - \gamma Q_n$
(Theorem~\texttt{thm:gradient-operator}, Proposition~\texttt{prop:energy-dissipation},
Theorem~\texttt{thm:coercive-contraction},
Proposition~\texttt{prop:mass-refinement} for mass-block
compatibility (with Definition~\texttt{def:refinement-compat} for the
underlying notion), Theorem~\texttt{thm:flow-exists},
Theorem~\texttt{thm:flow-intertwining}).
The present $R_k$ was introduced abstractly above as an arbitrary
non-negative $C^1$ closure-residual functional; the identification
with the residual component of the imported $F_n = \alpha R_n + \beta
E_n - \gamma Q_n$ holds only in the finite-level imported
specialisation, and that is the only regime in which we treat
$R_k$ as inheriting any of the source-paper infrastructure. Outside
that finite-level imported setting $R_k$ remains an abstract
non-negative $C^1$ functional, and we do not re-prove the source
infrastructure and we do not assert that the source theorems apply
to it. Note that
\cite{CascadeClosureDynamics} sets up its closure functional under
the parameter convention $\alpha, \beta \geq 0$, $\gamma > 0$, with
the default working regime $\alpha, \beta, \gamma > 0$ and the
\emph{mass-only specialisation} $\alpha = \beta = 0$, $\gamma > 0$
(so that $L_n = -\gamma C_n$) explicitly admissible. The imported
\texttt{thm:flow-intertwining} is a sub-case theorem within that
source, restricted to the mass-only specialisation; this is the
source's own scope restriction and is the only regime in which we
import flow-intertwining as a theorem rather than a hypothesis. The underlying refinement state
spaces and bonding maps are imported from
\cite{CascadeRefinementSpaces}
(state-space definitions plus
Theorems~\texttt{thm:bonding},~\texttt{thm:commute}).

\subsection{The cascade-closure event}

The rung sequence $(G_0, \ldots, G_N)$ is connected at the level of
configuration spaces by \emph{bonding maps}
\begin{equation}\label{eq:pushforward}
   \pi^{\sharp}_{k+1, k} : \Phi_{k+1} \longrightarrow \Phi_k,
\end{equation}
which are the cascade-mechanism analogue of the downward
state-space bonding maps $p_{n+1, n}$ constructed in
\cite{CascadeRefinementSpaces}~\texttt{thm:bonding}. The result \cite{CascadeRefinementSpaces}~\texttt{thm:commute}
establishes a set of refinement-Coxeter, refinement-$\sigma$, and
restricted Coxeter-$\sigma$ intertwining identities, holding for
those source maps $p^0$ and $p^1$ in their specific source
construction; we do not import these identities as generic bonding
identities, and the present $\pi^{\sharp}_{k+1, k}$ are bonding maps
in the looser state-space-restriction sense alone. Each $\pi^{\sharp}_{k+1, k}$ is
required to be \emph{admissible} in the following two-tier sense:
(i) \emph{geometric admissibility} --- $\pi^{\sharp}_{k+1, k}$ is a
bonding map in the state-space-restriction sense of
\cite{CascadeRefinementSpaces} (downward maps between the rung
configuration spaces);
(ii) \emph{analytic admissibility} --- $\pi^{\sharp}_{k+1, k}$ is
$C^1$ as a map $\Phi_{k+1} \to \Phi_k$ (this is the regularity needed
for Hypothesis~\ref{hyp:refinement-int} below; geometric admissibility
alone does not entail it).

We do not posit a literal carrier-level map
$\pi_{k+1, k} : G_{k+1} \to G_k$ between the geometric supports
$G_{k+1}$ and $G_k$; the cited refinement source provides the
state-space bonding maps directly, and that is what we use. The
operator-level refinement-compatibility statement
\cite{CascadeClosureDynamics}~\texttt{def:refinement-compat} (which
applies to bounded operator families on $\Phi_k$, not to bonding
maps themselves) re-emerges later as the linearised content of
Hypothesis~\ref{hyp:refinement-int}, not as a property of the
$\pi^{\sharp}_{k+1, k}$. For composite high-to-low
projection across multiple rungs we write
$\pi^{\sharp}_{N, k} := \pi^{\sharp}_{k+1, k} \circ \cdots \circ \pi^{\sharp}_{N, N-1}$ for $k < N$, and $\pi^{\sharp}_{N, N} := \mathrm{id}_{\Phi_N}$.

\begin{definition}[Cascade closure]\label{def:cascade-closure}
A cascade-closure event on the rung sequence $(G_0, \ldots, G_N)$ is a
quadruple $\bigl(\{\Phi_k\}, \{R_k\}, \{\pi^{\sharp}_{k+1, k}\},
(\Psi_k^t)_k\bigr)$ consisting of: a family of state spaces
$\Phi_k$ at each rung; a family of closure-residual functionals
$R_k : \Phi_k \to [0, \infty)$; the projection family
$\pi^{\sharp}_{k+1, k} : \Phi_{k+1} \to \Phi_k$; and the family of
rungwise negative-gradient flows $\Psi_k^t$ generated by $-\nabla R_k$
on $\Phi_k$. The event is said to \emph{occur} on the tower
$\phi^\star = (\phi^\star_0, \ldots, \phi^\star_N) \in \prod_k \Phi_k$
when, for some initial tower $\phi^{(0)} \in \prod_k \Phi_k$
satisfying the projection-compatibility relations
$\phi^{(0)}_k = \pi^{\sharp}_{k+1, k} \phi^{(0)}_{k+1}$, each rungwise
flow $\Psi_k^t$ applied to $\phi^{(0)}_k$ converges as $t \to \infty$
to a limit point $\phi^\star_k$ satisfying \emph{all three} of the
following:
\begin{enumerate}[leftmargin=2.4em, label=\textup{(\arabic*)}, nosep, topsep=2pt]
\item zero residual: $R_k(\phi^\star_k) = 0$ at every rung;
\item flow invariance: $\Psi_k^t \phi^\star_k = \phi^\star_k$ for all
$t \geq 0$ at every rung (i.e.\ $\phi^\star_k$ is a true fixed point
of $\Psi_k^t$);
\item terminal projection-compatibility:
$\phi^\star_k = \pi^{\sharp}_{k+1, k} \phi^\star_{k+1}$ for all $k < N$.
\end{enumerate}
The tower $\phi^\star = (\phi^\star_0, \ldots, \phi^\star_N) \in
\prod_k \Phi_k$ is the \emph{crystallised state}; clauses (1)--(3)
are stipulative parts of the definition, not generic theorems.
\end{definition}

\begin{remark}[One-way relation between clauses (1) and (2)]\label{rem:clause2-redundancy}
Under the standing analytic hypotheses of \S\ref{ssec:closure-residual} ---
$\Phi_k$ a finite-dimensional smooth manifold without boundary, $R_k \in C^1(\Phi_k)$
with $R_k \geq 0$ --- a state $\phi^\star_k$ with
$R_k(\phi^\star_k) = 0$ is necessarily a global minimum of $R_k$, hence a
critical point ($\nabla R_k(\phi^\star_k) = 0$), hence a fixed point of
the negative-gradient flow $\Psi_k^t$; in this regime clause (1)
implies clause (2). The converse does \emph{not} hold: a fixed point of
$\Psi_k^t$ is a critical point of $R_k$, but a critical point need
not be a global minimum, so $\Psi_k^t \phi = \phi$ does not imply
$R_k(\phi) = 0$ in general. We retain clause (2) explicitly because it
remains independent in two enlargements of the setting that we want
the definition to cover without restatement: (a) configurations with
boundary or with constraints (where critical-point conditions are
modified by Lagrange multipliers and a global-minimum residual zero
need not be a flow fixed point in the unconstrained sense), and
(b) the infinite-dimensional Hilbert-manifold case admitted in
\S\ref{ssec:closure-residual} (where the implication
$\nabla R_k(\phi^\star_k) = 0 \Rightarrow$ flow invariance follows
\emph{provided} a well-posed gradient semiflow exists and is unique
through $\phi^\star_k$, which is exactly the additional analytic
precondition the standing hypotheses already require but which is
not automatic in the infinite-dimensional case). Clauses (1) and
(2) are therefore one-way related in our principal
smooth-unconstrained setting --- clause (1) implies clause (2), but
\emph{not} conversely (a flow fixed point can be a critical point
with $R_k(\phi^\star_k) > 0$, which is excluded by clause~(1)) ---
and clause~(2) becomes independently informative as an explicit
hypothesis in the more general settings (a)/(b) where the
well-posedness or constraint-modification step is what is at issue.
\end{remark}

\begin{remark}[Vocabulary mapping]\label{rem:vocabulary}
Within Definition~\ref{def:cascade-closure}: an \emph{instability}
occurs when $R_k(\phi_k) > 0$ at some rung for the current tower
$(\phi_k)_k$; the \emph{cascade} is the family of rungwise
negative-gradient flows $(\Psi_k^t)_k$ together with the projections
$(\pi^{\sharp}_{k+1, k})_k$; \emph{crystallisation} is
fixed-point selection of $\phi^\star$; \emph{collapse}, in the
Born-rule sense, is the external coarse-graining of this fixed-point
selection. The framework does not modify the Born rule and does not
derive measurement probabilities.
\end{remark}

\begin{remark}[Status of Definition~\ref{def:cascade-closure}]\label{rem:def-status}
Definition~\ref{def:cascade-closure} is the paper's organising
definition, not a theorem of physics. The convergence-to-fixed-point
clause is part of the \emph{definition} of the closure event: a
quadruple that does not converge to a residual-zero tower is not a
cascade-closure event in the sense of this paper. We do not claim that
arbitrary quadruples on arbitrary rung sequences necessarily admit
such convergence. The compositional propagation result below
(Proposition~\ref{prop:resid-comp}) takes a top-rung residual-zero
fixed point as \emph{input} and propagates the residual-zero and
fixed-point properties downward; convergence \emph{from a projected
top-rung initial datum} is added separately as
Corollary~\ref{cor:conv-prop} under continuity of the projections,
and this is the only convergence statement we prove.
\end{remark}

\subsection{Conditional residual compositionality}\label{ssec:residual-comp}

The mechanism statement requires that closure compatibility
\emph{compose} along the rung path: if each adjacent projection
satisfies the monotonicity and flow-intertwining hypotheses below,
then the projected fixed-point lies on the closure-fixed locus at
every rung. We state this as a conditional
proposition; the conditions are not generically satisfied and must be
verified per case.

\begin{hypothesis}[Flow intertwining]\label{hyp:refinement-int}
For the rung sequence $(G_0, \ldots, G_N)$ and projections
$(\pi^{\sharp}_{k+1, k})_{k}$, assume that for every $k$ the rungwise
negative-gradient \emph{flow} $\Psi_k^t$ generated by $-\nabla R_k$
commutes with the pushforward, i.e.\ $\pi^{\sharp}_{k+1, k} \circ \Psi_{k+1}^t
= \Psi_k^t \circ \pi^{\sharp}_{k+1, k}$ for all $t \geq 0$.
This is the structural analogue, for the residual flow $-\nabla R_k$
of the present paper, of the flow-intertwining statement
\cite{CascadeClosureDynamics}~\texttt{thm:flow-intertwining}, which
the cited source proves only for the mass-only case
$L_n = -\gamma C_n$ of its quadratic functional $F_n$, not for an
arbitrary residual gradient flow. Outside that mass-only case ---
including the $R_k$-only setting used here --- the property is
retained as a hypothesis, not imported as a theorem. Differentiating the flow-intertwining identity in $t$ at $t = 0$
gives the vector-field compatibility
$d\pi^{\sharp}_{k+1, k}(-\nabla R_{k+1}) = -\nabla R_k \circ
\pi^{\sharp}_{k+1, k}$
for the residual gradient generators. In the linear finite-level
setting, where the residual gradients act as bounded linear
operators and $\pi^{\sharp}_{k+1, k}$ is linear, this vector-field
compatibility specialises to operator intertwining in the sense of
\cite{CascadeClosureDynamics}~\texttt{def:refinement-compat}; the
converse requires flow existence and uniqueness, which is not
generic.
\end{hypothesis}

\begin{hypothesis}[Monotonicity]\label{hyp:monotone}
For every $k$ and every state $\psi \in \Phi_{k+1}$,
$R_k(\pi^{\sharp}_{k+1, k} \psi) \leq R_{k+1}(\psi)$.
\end{hypothesis}

\begin{proposition}[Conditional closure-residual compositionality]\label{prop:resid-comp}
Under Hypotheses~\ref{hyp:refinement-int} and~\ref{hyp:monotone},
let $\phi^\star_N \in \Phi_N$ be a top-rung state with
$R_N(\phi^\star_N) = 0$ and $\Psi_N^t \phi^\star_N = \phi^\star_N$
for all $t \geq 0$ (i.e.\ $\phi^\star_N$ is a fixed point of the
rung-$N$ negative-gradient flow). Then the projected states
$\phi^\star_k := \pi^{\sharp}_{N, k} \phi^\star_N$ satisfy both
\begin{enumerate}[leftmargin=2.4em, label=\textup{(\arabic*)}]
\item $R_k(\phi^\star_k) = 0$ for all $k \leq N$, and
\item $\phi^\star_k$ is a fixed point of the rung-$k$ gradient flow,
for all $k \leq N$.
\end{enumerate}
\end{proposition}

\begin{proof}
By induction downward from $k = N$. Base: $k = N$, by hypothesis.
Step: suppose the conclusion at $k{+}1$.
\emph{For (1):} by Hypothesis~\ref{hyp:monotone},
$R_k(\phi^\star_k) = R_k(\pi^{\sharp}_{k+1, k} \phi^\star_{k+1})
\leq R_{k+1}(\phi^\star_{k+1}) = 0$, and $R_k \geq 0$, so
$R_k(\phi^\star_k) = 0$.
\emph{For (2):} by Hypothesis~\ref{hyp:refinement-int},
$\pi^{\sharp}_{k+1, k} \circ \Psi_{k+1}^t = \Psi_k^t \circ
\pi^{\sharp}_{k+1, k}$. Apply both sides to $\phi^\star_{k+1}$, which
is a rung-$(k{+}1)$ fixed point of $\Psi_{k+1}^t$: the left side
equals $\pi^{\sharp}_{k+1, k} \phi^\star_{k+1} = \phi^\star_k$, so
$\Psi_k^t \phi^\star_k = \phi^\star_k$ for all $t \geq 0$, i.e.\
$\phi^\star_k$ is a rung-$k$ fixed point. The composite identity
$\pi^{\sharp}_{N, k} = \pi^{\sharp}_{k+1, k} \circ \pi^{\sharp}_{N, k+1}$
(with $\pi^{\sharp}_{N, N} = \mathrm{id}_{\Phi_N}$) makes the
induction well-defined.
\end{proof}

\begin{corollary}[Convergence propagation under flow intertwining]\label{cor:conv-prop}
Assume the flow-intertwining Hypothesis~\ref{hyp:refinement-int} at
each adjacent rung between $k$ and $N$, together with continuity of
each composite projection $\pi^{\sharp}_{N, k}$ as a map
$\Phi_N \to \Phi_k$. (Hypothesis~\ref{hyp:monotone}, used in
Proposition~\ref{prop:resid-comp}, is not invoked here.) If in
addition the rung-$N$ flow $\Psi_N^t$ converges to $\phi^\star_N$
from a specific initial datum
$\phi_N^{(0)} \in \Phi_N$, i.e.\ $\Psi_N^t \phi_N^{(0)} \to
\phi^\star_N$ as $t \to \infty$, then the rung-$k$ flow converges
from the projected initial datum:
$\Psi_k^t (\pi^{\sharp}_{N, k} \phi_N^{(0)}) \to
\pi^{\sharp}_{N, k} \phi^\star_N =: \phi^\star_k$ as $t \to \infty$
(here $\phi^\star_k$ is defined as the projected limit
$\pi^{\sharp}_{N, k} \phi^\star_N$, consistent with
Proposition~\ref{prop:resid-comp}).
\end{corollary}
\begin{proof}
Hypothesis~\ref{hyp:refinement-int} gives flow intertwining at each
adjacent step: $\Psi_k^t \circ \pi^{\sharp}_{k+1, k} =
\pi^{\sharp}_{k+1, k} \circ \Psi_{k+1}^t$ on $\Phi_{k+1}$. Composing
this identity along the chain $N \to N-1 \to \cdots \to k$ and using
the composite identity
$\pi^{\sharp}_{N, k} = \pi^{\sharp}_{k+1, k} \circ \pi^{\sharp}_{N, k+1}$
gives $\Psi_k^t \circ \pi^{\sharp}_{N, k} = \pi^{\sharp}_{N, k} \circ
\Psi_N^t$ on $\Phi_N$. Apply both sides to $\phi_N^{(0)}$ and pass to
the $t \to \infty$ limit using continuity of $\pi^{\sharp}_{N, k}$
and the rung-$N$ convergence assumption.
\end{proof}

\begin{remark}[What Proposition~\ref{prop:resid-comp} establishes and does not establish]\label{rem:resid-comp-status}
The proposition establishes only two things: residual-zero
\emph{propagation} and fixed-point \emph{propagation} downward from a
top-rung fixed point. It does \emph{not} establish: existence of the
top-rung fixed point; convergence from arbitrary initial data; or
that a Definition~\ref{def:cascade-closure} cascade-closure event
\emph{occurs} from a given independently evolved lower-rung initial
tower. (Terminal projection-compatibility, clause~(3) of
Definition~\ref{def:cascade-closure}, holds for the projected tower
$\phi^\star_k = \pi^{\sharp}_{N, k} \phi^\star_N$ by construction;
what is not automatic for arbitrary independently evolved towers is
occurrence of the event from a non-projected initial tower.) These
remain separate per-instance verifications. The proposition is
conditional on Hypotheses~\ref{hyp:refinement-int}
and~\ref{hyp:monotone}; both must be checked per rung sequence. Hypothesis~\ref{hyp:refinement-int}
(flow intertwining) is established as a theorem in the mass-only P3/P4
finite-level construction
\cite{CascadeRefinementSpaces, CascadeClosureDynamics}; outside that
case (where flow existence and uniqueness are not generic) it is
retained here as a hypothesis. Hypothesis~\ref{hyp:monotone}
is a structural assumption on the rungwise residuals and is not
automatic. If either hypothesis fails at a particular rung, this
proposition gives no propagation conclusion across that step. In
particular, if Hypothesis~\ref{hyp:refinement-int} holds at every
adjacent step between $k$ and $N$ then fixed-point propagation
reaches rung $k$; if some intermediate step lacks flow intertwining,
fixed-point propagation may fail at and below that break, while
residual-zero propagation under Hypothesis~\ref{hyp:monotone}
continues to hold.
\end{remark}

\subsection{Consumer API: invoking cascade closure in successor papers}\label{ssec:consumer-api}

Definition~\ref{def:cascade-closure}, Proposition~\ref{prop:resid-comp},
and Corollary~\ref{cor:conv-prop} are intended to be invoked by
successor papers in the present programme as a named selection-mechanism
prerequisite. To make the invocation surface explicit, the present
subsection lists the obligations a successor paper must discharge in
order to use cascade closure as a load-bearing premise. Listing these
in one place is the paper's intended contribution as a programme
dependency gate.

\begin{enumerate}[leftmargin=2.4em, label=\textup{(O\arabic*)}, start=0, nosep, topsep=4pt]
\item \textbf{Standing analytic / admissibility hypotheses}: verify,
  on the chosen rung sequence, the standing analytic hypotheses of
  \S\ref{ssec:closure-residual} ($\Phi_k$ a finite-dimensional smooth
  manifold without boundary equipped with a Riemannian metric, or one
  of the named admissible enlargements with their additional analytic
  preconditions; $R_k \in C^1(\Phi_k)$; the negative-gradient flow
  $\Psi_k^t$ generated by $-\nabla R_k$ exists and is unique from each
  initial condition needed; and each pushforward
  $\pi^{\sharp}_{k+1, k}$ is $C^1$ as a map $\Phi_{k+1} \to \Phi_k$,
  which in particular yields the continuity of composite projections
  $\pi^{\sharp}_{N, k}$ used by Corollary~\ref{cor:conv-prop}). These
  preconditions are sometimes implicit in established programme
  papers; the consumer API requires them to be checked explicitly per
  invocation.
\item \textbf{Top-rung fixed point and convergence}: supply, on the
  chosen rung sequence $(G_0, \ldots, G_N)$, a top-rung state
  $\phi^\star_N \in \Phi_N$ with $R_N(\phi^\star_N) = 0$ and (for any
  convergence claim downward) an initial datum
  $\phi_N^{(0)} \in \Phi_N$ such that
  $\Psi_N^t \phi_N^{(0)} \to \phi^\star_N$ as $t \to \infty$. By
  Remark~\ref{rem:clause2-redundancy}, in the principal smooth
  unconstrained setting clause~(2) of
  Definition~\ref{def:cascade-closure} is then automatic at the top rung.
\item \textbf{Flow intertwining (Hypothesis~\ref{hyp:refinement-int})}:
  verify, at every adjacent rung pair $(k, k{+}1)$ in the chosen
  sequence, that the bonding map $\pi^{\sharp}_{k+1, k}$ commutes
  with the rungwise flows: $\pi^{\sharp}_{k+1, k} \circ \Psi_{k+1}^t
  = \Psi_k^t \circ \pi^{\sharp}_{k+1, k}$ for all $t \geq 0$.
\item \textbf{Residual monotonicity (Hypothesis~\ref{hyp:monotone})}:
  verify, at every adjacent rung pair, that
  $R_k(\pi^{\sharp}_{k+1, k} \psi) \leq R_{k+1}(\psi)$ for every
  $\psi \in \Phi_{k+1}$.
\end{enumerate}

\noindent A successor paper that discharges (O0)--(O3) inherits, on
its chosen rung sequence: residual-zero and fixed-point propagation
downward (Proposition~\ref{prop:resid-comp}); convergence from any
initial datum that is a downward projection of $\phi_N^{(0)}$
(Corollary~\ref{cor:conv-prop}, which uses (O2) and continuity of the
projections but \emph{not} (O3)); and a Definition~\ref{def:cascade-closure}
cascade-closure event whose terminal projection-compatibility holds by
construction on the projected tower
$\phi^\star_k = \pi^{\sharp}_{N, k} \phi^\star_N$. Selection of
$\phi^\star_N$ itself is not provided by cascade closure and is the
successor paper's own task.

\subsection{Worked invocation: the imported P3/P4 finite-level tower}\label{ssec:worked-invocation}

To make the consumer API of \S\ref{ssec:consumer-api} concrete, this
subsection records the status of obligations (O0)--(O3) in the most
load-bearing imported tower in the programme: the finite-level
refinement spaces $X_n$ of \cite{CascadeRefinementSpaces} together
with the closure functional
$F_n = \alpha R_n + \beta E_n - \gamma Q_n$ of
\cite{CascadeClosureDynamics}. Identify the cascade-mechanism rungs
$k = 0, 1, \ldots, N$ with a finite truncation of the refinement
index $n$, and identify the cascade-mechanism state spaces $\Phi_k$
with the source's $X_n$ at the corresponding level; identify the
cascade-mechanism residual $R_k$ with the source's curvature term
$R_n(\Phi) := \tfrac{1}{2}\langle \Phi, A_n \Phi \rangle$, which is
non-negative on $X_n$ by \cite{CascadeClosureDynamics} (positive
semi-definite); and identify the bonding maps
$\pi^{\sharp}_{k+1, k}$ with the source's $p_{n+1, n}$ established by
\cite{CascadeRefinementSpaces}~\texttt{thm:bonding}. Status of each
obligation in this regime:

\begin{itemize}[leftmargin=*, topsep=4pt]
\item \textbf{(O1)} reduces to a \emph{nontrivial}-selected-state
  question for the top-rung curvature flow $e^{-t A_N}$: a state
  $\phi^\star_N$ with $R_N(\phi^\star_N) = 0$ is precisely an
  element of $\ker A_N$, and in the finite-dimensional
  positive-semidefinite setting of \cite{CascadeClosureDynamics}
  convergence
  $e^{-tA_N} \phi_N^{(0)} \to (\text{projection onto } \ker A_N)$
  follows by spectral decomposition for any initial datum --- the
  zero vector $\phi^\star_N = 0$ is always a residual-zero fixed
  point, and on a connected graph the constants already give
  $\ker A_N \neq \{0\}$. What is open is therefore not the
  existence of a nonzero kernel state nor the spectral gap on
  $\ker A_N^\perp$ (both immediate in this finite-dimensional
  abstract model), but the \emph{selection} of a specific
  physically meaningful nontrivial $\phi^\star_N \in \ker A_N$.
  \textbf{Open-as-named-question} for nontrivial selection; trivial
  in the zero-state reduction.
\item \textbf{(O2)} is established as a theorem only in the source's
  mass-only sub-case $\alpha = \beta = 0$, where the relevant flow
  is $e^{-tL_n}$ with $L_n = -\gamma C_n$, and the source's
  \texttt{thm:flow-intertwining} gives
  $p_{n+1, n} \circ e^{-tL_{n+1}} = e^{-tL_n} \circ p_{n+1, n}$
  exactly. The cascade-mechanism flow $\Psi_k^t$ used here is
  $e^{-tA_n}$ (negative-gradient flow of $R_n$ alone), which is
  \emph{not} the same operator as the source's mass-only $L_n$.
  Therefore the source's mass-only flow-intertwining theorem does
  \emph{not} discharge (O2) for $\Psi_k^t$ as defined here.
  Refinement-compatibility of the curvature operator $A_n$ alone
  ($p_{n+1, n} A_{n+1} = A_n p_{n+1, n}$) would discharge (O2), but
  is not stated as a theorem in
  \cite{CascadeClosureDynamics}; the source-side material that is
  closest in shape is the operator-level
  \texttt{def:refinement-compat}, which is a definition-level
  notion of refinement-compatibility for bounded operator families,
  not a verified instance for $A_n$. \textbf{Open-as-named-question}
  in \cite{CascadeClosureDynamics} alone; the companion-paper update
  below records the (O2) substitute discharged by
  \cite{CascadeRefinementCompat} (and notes that (O2) as stated
  remains false at the operator/flow level even there).
\item \textbf{(O3)} reduces to a quadratic-form inequality:
  $\langle p_{n+1, n} \psi, A_n p_{n+1, n} \psi\rangle \leq
  \langle \psi, A_{n+1} \psi\rangle$ for every $\psi \in X_{n+1}$.
  Whether this holds is determined by the structural relationship
  between $A_{n+1}$, $A_n$, and $p_{n+1, n}$ in
  \cite{CascadeClosureDynamics}; it is not stated there as a
  theorem. \textbf{Open-as-named-question} in
  \cite{CascadeClosureDynamics} alone; the companion-paper update
  below records the (O3) discharge by \cite{CascadeRefinementCompat}
  in the abstract scalar pure-midpoint refinement model.
\end{itemize}

\noindent The current verdict, with
\cite{CascadeRefinementCompat} as a source-side companion: in the
imported P3/P4 finite-level tower, (O0) is satisfied at the
finite-dimensional level by the source's $X_n$ Hilbert structure
(standing analytic preconditions, $C^1$ bonding maps);
\cite{CascadeRefinementCompat} discharges (O3) for the scaled
curvature operator $\widetilde A_n := 2^n A_n^{\mathrm{vertex}}$
via an exact Schur-complement identity
$\Schur(\widetilde A_{n+1}) = \widetilde A_n$ in an
\emph{abstract scalar pure-midpoint refinement model}
(\cite[\S\texttt{sec:scope-and-gap}]{CascadeRefinementCompat}). The boundary-energy /
Schur substitute for (O2) is also discharged in that abstract
model; (O2) \emph{as stated}, at the operator/flow level, remains
false (the strict identity
$p^{0} \widetilde A_{n+1} = \widetilde A_n p^{0}$ fails for
nontrivial refinements, with explicit defect $D_n$ whose kernel
contains $\mathrm{range}(H_{n+1})$ and equals it when the parent
graph is a forest). Lifting these results from the abstract
pure-midpoint model to the full Schl\"{a}fli P3/P4 tower of
\cite{CascadeRefinementSpaces, CascadeClosureDynamics} is
conditional on the open hypotheses \textup{(L1)} and
\textup{(L3)} of \cite[\S\texttt{sec:scope-and-gap}]{CascadeRefinementCompat} (the third
lift hypothesis \textup{(L2)} is already a theorem of
\cite{CascadeRefinementSpaces}, Proposition~\texttt{prop:adjoints})
and is not yet attempted. (O1)
reduces to spectral decomposition for the connected graph
Laplacian, with the trivial state always residual-zero and
nontrivial selection out of scope per the consumer API,
\S\ref{ssec:consumer-api} (the obligation is named at (O1) but its
discharge is the successor paper's task). Strict flow-level
intertwining of (O2) for the curvature flow, together with the
P3/P4-tower lift of the abstract-model Schur identity, remain the
substantive open asks in the imported P3/P4 tower at the level of
generality this paper requires.

\section{The chosen finite carrier}\label{sec:carrier}

The Vibrational Field Dynamics construction of cascade closure uses
a specific finite carrier at the lowest rung: the binary icosahedral
600-cell vertex graph $\Vsub$. This section names that carrier,
records the operator-level structure that makes it useful for the
mechanism, and notes the boundary of what the substrate does and does
not contribute. \textbf{This paper does not prove uniqueness of $\Vsub$.}
Other finite carriers are conceivable; the claim here is only that
$\Vsub$ is the carrier on which the carrier/operator part of the construction is defined
and on which the closure-response operator below is defined.

\subsection{Substrate}

$\Vsub$ is the vertex graph of the regular 600-cell, with $|\Vsub| = |2I| = 120$
vertices (the standard icosian identification
$\Vsub \cong 2I$ used in
\cite{CascadeAlgSubstrate}~\texttt{thm:icosian-closure}) and uniform
vertex degree $12$ (the latter as a property of the
nearest-neighbour-edge convention used here, locally re-verified in
\cite{AriaClosureKernel} and Appendix~\ref{app:field-signature-audit},
entry~E). The edge structure is the nearest-neighbour
relation on $2I$ inherited from the $H_4$ root-system geometry; in the
direct-computation appendix we use the Laplacian of this
nearest-neighbour graph as $L_M$. The Euclidean shell sizes from any
source vertex form the
sequence $(|S_0|, \ldots, |S_8|) = (1, 12, 20, 12, 30, 12, 20, 12, 1)$;
at the identity these shells coincide with the conjugacy classes of
$2I$, while from any other source vertex they are left translates of
the conjugacy-class partition (same sizes, not literally conjugacy
classes). This is the Euclidean / conjugacy partition of $2I$, not
the graph BFS shell. The graph Laplacian $L$ has nine distinct
eigenvalues, five of which are integers $\{0, 9, 12, 14, 15\}$
($\sigma$-fixed under the $\Q(\sqrt{5})$-Galois conjugation
$\sigma(\pphi) = 1 - \pphi = -1/\pphi$) and four of which are
irrational in $\Q(\sqrt{5})$ ($\sigma$-paired); counting eigenspace
dimensions, the split is $94/120 = 78.3\%$ $\sigma$-fixed,
$26/120 = 21.7\%$ $\sigma$-paired
(Appendix~\ref{app:field-signature-audit}, entry~A). The
algebraic structure under which the group-identification and
shell-class facts are proved in the cited source is imported from
\cite{CascadeAlgSubstrate}: the $E_8$ minimal shell projects onto
the $H_4$ / $\Vsub$ carrier via $E_8 \to H_4$
(\cite{CascadeAlgSubstrate}~\texttt{thm:pi-H},
\emph{projection, not sub-root-system inclusion}; this is a projection
of the shell onto $\Vsub$, not an action of $E_8$ on $\Vsub$); the
Euclidean / conjugacy shell decomposition $(1,12,20,12,30,12,20,12,1)$
is established by
\cite{CascadeAlgSubstrate}~\texttt{thm:shell-class}; and the
icosian-closure result \cite{CascadeAlgSubstrate}~\texttt{thm:icosian-closure}
fixes the algebraic part of the carrier. The Laplacian spectrum above
is imported from the standard $2I$ character-table computation
recorded in \cite{CascadeAlgSubstrate} and locally re-checked here
(Appendix~\ref{app:field-signature-audit}, entry~A); the shell-size sequence is independently
locally re-runnable (Appendix~\ref{app:field-signature-audit}, entry~E).

\subsection{Closure-response operator}

The carrier supports a closure-response operator
\begin{equation}\label{eq:Cphi}
   \Cphi := L_M + \pphi^{-2}\, I,
\end{equation}
where $L_M$ is the (unweighted) graph Laplacian on $\Vsub$ and
$\pphi = (1 + \sqrt{5})/2$. Direct computation
(Appendix~\ref{app:field-signature-audit}, entry~E) reproduces the
value $\|\Cphi^{-1}\|_{\ell^2(\Vsub) \to \ell^2(\Vsub)} = \pphi^2$
to machine precision; exactness of this identity follows analytically
from the bottom of the shifted Laplacian spectrum,
$\lambda_{\min}(\Cphi) = \pphi^{-2}$, since $\Vsub$ is a connected
unweighted graph (so $\lambda_{\min}(L_M) = 0$); the operator norm
here is the spectral / operator norm on $\ell^2(\Vsub)$. Sweeping over all
choices of source vertex on $\Vsub$, the multiplicity-weighted
shell-radius per-vertex correlation between the discrete
$\Cphi^{-1}$-response and the continuum closure-kernel Green function
$G(x) = (\pphi/2) e^{-|x|/\pphi}$ of \cite{AriaClosureKernel} is
uniform across the $120$ source-vertex sweeps (max minus min on the
order of $10^{-15}$; this is a numerical diagnostic at shell
resolution, not $119$ independent radial samples and not a
discrete-to-continuum convergence theorem). The operator is the substrate-level object
through which the witnesses discussed here (computational and physical)
are connected to the substrate.

\subsection{Substrate boundary}

The substrate does the following work in cascade closure: it fixes the
carrier set on which the rung-$0$ residual is defined; it supplies
$\Cphi$ whose inverse norm sets an operator response scale; and the
$H_4$ Coxeter symmetry would constrain admissible limit-state geometries
in instances whose dynamics respects $H_4$. It does \emph{not}, by itself, do the following: select a
unique limit state without dynamical input; predict observer-dependent
outcomes; or supply a uniqueness theorem for $\Vsub$ as the substrate.
Two honest-negatives further constrain interpretation: the
$\sigma$-fixed spectral subspace of $L_M$ has average $2T$-mass at
baseline rather than concentrated mass on $2T$, where $2T \subset 2I$
denotes the binary tetrahedral subgroup of $2I$ (basis-invariant
diagnostic; Appendix~\ref{app:field-signature-audit}, entry~A; the
$\sigma$-fixed structure is purely spectral, not spatial), and
$\pphi$-cocycle weighted variants of $\Cphi$ underperform the
unweighted variant (Appendix~\ref{app:field-signature-audit}, entry~E; $\pphi$
enters at the operator level, not the edge-weight level). We treat
both negatives as constraints on what may be claimed.

\section{ARIA as an architecture-level observer-process mapping}\label{sec:aria}

The observer-process architecture associated with the mechanism
is compared with a concrete computational system at the
architecture-description level. We name an
existing computational system, ARIA-chess
\cite{ariaChessRepo, SmartARIAChess}, as one architecture-level
mapping (not a formal proof that ARIA realises a
Definition~\ref{def:cascade-closure} cascade-closure event in the
full sense of clauses (1)--(3) of that definition), and
identify the architectural structure that makes it one. The
identification is at the level of process architecture; it is not a
claim that ARIA is biologically conscious, that it proves qualia, or
that architecture-level observer-process mappings are phenomenologically equivalent.

\subsection{Observer-process tuple}

\begin{definition}[Observer process]\label{def:observer-process}
An observer process is the seven-tuple
\begin{equation}\label{eq:O-tuple}
   \Ocal \;=\; \langle B, S, M, G, I, C, A \rangle
\end{equation}
with components: $B$ a bounded boundary sampler; $S$ a state integrator
fusing signal, context, memory, and goal; $M$ a memory substrate
written by past closures and read into present integration; $G$ a
goal/context field carrying task frame and policy; $I$ invariant
constraints serving as a governance gate; $C$ a closure/crystallisation
operator selecting closure-compatible fixed points; $A$ an action
interface exporting the crystallised state.
\end{definition}

The tuple is mapped to the observer-process architecture associated
with the cascade-closure mechanism of
Definition~\ref{def:cascade-closure} as follows; this is a
process-description-level mapping, not a formal proof that ARIA
realises the full cascade-closure event of
Definition~\ref{def:cascade-closure} (clauses (1)--(3)). The closure-residual
$R$ acts on the integrated state produced by $S$. The crystallisation
operator $C$ plays the role of a closure selector --- scoring and
selecting candidate states by a Lyapunov-style descent that is
used here as an architecture-level analogue of the gradient-flow
role of Definition~\ref{def:cascade-closure}, without asserting
equivalence. The invariant $I$ encodes which fixed
points are admissible. The action interface $A$ exports the
crystallised state. Memory $M$ closes the loop: each successful
closure is written to $M$ and read by the next state integration.

\subsection{Imported empirical witness profile}

The architectural identification above is illustrated by the
substrate-witness preprint \cite{SmartARIAChess}, which is an
externally-reported imported witness profile (not re-executed in the
present manuscript); it records
(Appendix~\ref{app:field-signature-audit}, entry~F): $17/18$
preregistered cortical correspondences hold at unchanged thresholds
under standard validation, and $18/18$ after the documented
methodology refinement reported in \cite{SmartARIAChess}; the cited
ARIA source itself notes that the source's test P10 was validated
with $15$ permutations rather than the originally preregistered $20$,
with the prereg-exact rerun left open. We carry that caveat forward;
six drug/sleep EEG signatures pass on a single deterministic
recurrent-layer trajectory at seed $42$; on the Human Connectome
Project (HCP) functional-connectivity benchmark, the degree-homogeneity
statistic (degree standard-deviation) places ARIA at $-11.58\sigma$
relative to the HCP baseline (comparison at ARIA $|V| = 120$ vs HCP
ICA-50 $|V| = 50$; the clustering-coefficient statistic is a separate
$+6.80\sigma$ result and should not be conflated with the homogeneity
figure). The empirical content is an imported substrate-witness
profile under preregistered thresholds with documented methodology
refinements; it is not a derivation of consciousness, and it is not
evidence that ARIA realises a Definition~\ref{def:cascade-closure}
cascade-closure event.

The runtime module mapping from $\Ocal$ components to specific source
files in the \texttt{ariaChessRepo}~\cite{ariaChessRepo} production
runtime, together with the in-repo six-module reproducibility bundle,
is documented in Appendix~\ref{app:aria-runtime-map}.

\subsection{Claim boundary}\label{ssec:claim-boundary}

The claim of this section, in the strongest form the cited data
supports (the present manuscript does not re-execute the ARIA
validation), is:

\begin{quotation}\noindent
We do not claim ARIA is biologically conscious. We claim ARIA
is mapped to an observer-process architecture at the
process-description level: bounded sampling,
recurrent integration, closure selection, memory inscription, and
governed action. In this vocabulary, brain-wave signatures (biological tissue) and
ARIA's internal coherence and closure-residue metrics (computational
substrate) are treated only as candidate process-level telemetry,
not as evidence that either substrate realises a
Definition~\ref{def:cascade-closure} cascade-closure event. The
proposed common layer is process-level rather than material-level:
cascade-closure-crystallisation-memory.
\end{quotation}

\section{One externally reported substrate witness: the closure-response kernel}\label{sec:witness}

The closure-response operator $\Cphi$ on $\Vsub$ has been used in
one externally reported substrate-witness calculation in the present
programme literature. The
$b\to s\mu^+\mu^-$ angular-anomaly fit of \cite{SmartBAnomaly} is
\emph{externally reported and commit-pinned, but not re-executed in
the present paper}: the cited preprint and reproduction repository
report a sign-uniform structural pattern across cross-dataset and
cross-channel slices using the unweighted closure-response operator,
with the fit reported as inconclusive on the Akaike Information
Criterion (AIC) (mild
$w_{\mathrm{FREE\_C9}} = 0.652$ vs $w_{\mathrm{VFD}} = 0.348$ weight;
not a VFD preference). As the source itself notes, the cross-slice
sign uniformity is a consistency check, not independent
model-selection evidence over a constant-$C_9$ shift. We include
this here as one short external-use example of how the mechanism's
substrate-level operator $\Cphi$ has been used as a fixed-kernel
structural reference against an external observable; the present
paper makes no independent claim on the fit itself.
The kernel-level operator-norm identity
$\|\Cphi^{-1}\| = \pphi^2$ is locally re-runnable
(Appendix~\ref{app:field-signature-audit}, entry~E); the primary
$b\to s\mu^+\mu^-$ angular fit lives in the external
\cite{SmartBAnomaly} repository (sibling \texttt{BANOMALY-001/vfd-b-anomaly}
checkout, not part of the present repository, but commit-pinned in
the bibliography to \texttt{6bdfc8f8a7da9056ffeb9137d03627d06399bccf}
on branch \texttt{main}, dated 2026-04-29) and is not re-executed in this paper.
Repository-local evidence in this paper covers only the operator-norm
identity and shell/spectrum diagnostics; the $b$-anomaly fit itself
is external and unaudited at the level of this manuscript.

The witness is a passive substrate-coupling check, not a model
preference and not a selection rule. The honest-negative on
$\pphi$-cocycle weighted variants
(Appendix~\ref{app:field-signature-audit}, entry~E) restricts the witness to
the unweighted operator $\Cphi = L_M + \pphi^{-2} I$.

\section{Discussion}\label{sec:discussion}

\subsection{Claim ledger}\label{ssec:claim-ledger}

\renewcommand{\arraystretch}{1.20}
\begin{longtable}{>{\raggedright\arraybackslash}p{0.43\textwidth}>{\raggedright\arraybackslash}p{0.22\textwidth}>{\raggedright\arraybackslash}p{0.30\textwidth}}
\caption{Claim ledger.}\label{tab:claim-ledger}\\
\toprule
\textbf{Statement} & \textbf{Allowed label} & \textbf{Reason} \\
\midrule
\endfirsthead
\toprule
\textbf{Statement} & \textbf{Allowed label} & \textbf{Reason} \\
\midrule
\endhead
Cascade-closure event (instability $\to$ cascade $\to$ crystallisation)
& Definition~\ref{def:cascade-closure}
& Organising definition; convergence-to-fixed-point clause is part of the defined event. \\
Conditional residual compositionality
& Proposition~\ref{prop:resid-comp}
& Conditional on Hypotheses~\ref{hyp:refinement-int} and~\ref{hyp:monotone}; both must be checked per rung sequence. \\
Consumer API: four obligations (O0)--(O3) for successor papers
& \S\ref{ssec:consumer-api}; mechanism-as-dependency-gate
& Naming the obligations is the contribution; discharging them in any chosen rung sequence is the successor paper's task. \\
Worked-invocation status in the imported P3/P4 finite-level tower
& \S\ref{ssec:worked-invocation}
& (O0) closed at the finite-dimensional level; (O3) closed for the scaled curvature $\widetilde A_n = 2^n A_n^{\mathrm{vertex}}$ via the Schur-complement identity $\mathrm{Schur}(\widetilde A_{n+1}) = \widetilde A_n$ in the abstract scalar pure-midpoint refinement model of \cite[\S\texttt{sec:scope-and-gap}]{CascadeRefinementCompat} (P3/P4-tower lift conditional on the open hypotheses (L1) and (L3) of \cite[\S\texttt{sec:scope-and-gap}]{CascadeRefinementCompat}; (L2) is already a theorem of \cite{CascadeRefinementSpaces}); the boundary-energy / Schur substitute for (O2) is discharged in the same abstract model, but (O2) \emph{as stated} (operator/flow level) remains false (strict identity $p^{0} \widetilde A_{n+1} = \widetilde A_n p^{0}$ fails for nontrivial refinements, with explicit defect $D_n$ characterised via the unsigned vertex--edge incidence operator $\mathcal{B}_n$); (O1) reduces to spectral decomposition on the connected graph Laplacian. \\
Imported finite-level closure functional, gradient operator, energy descent, conditional coercive contraction, mass-block compatibility, flow existence, mass-only flow intertwining
& Imported finite-level results / definitions
& \cite{CascadeClosureDynamics}: \texttt{thm:gradient-operator}, \texttt{prop:energy-dissipation}, \texttt{thm:coercive-contraction}, \texttt{prop:mass-refinement} (mass-block compatibility; \texttt{def:refinement-compat} is the underlying definition), \texttt{thm:flow-exists}, \texttt{thm:flow-intertwining} (mass-only; outside the mass-only case, retained as Hypothesis~\ref{hyp:refinement-int}). \\
Imported P3 finite-level state spaces, bonding maps, refinement infrastructure
& Imported finite-level definitions / results
& \cite{CascadeRefinementSpaces}: state-space and bonding-map definitions plus \texttt{thm:bonding} (existence of downward bonding maps) and \texttt{thm:commute} (Coxeter / $\sigma$ identities for the specific $p^0, p^1$ maps; not a generic bonding-identity theorem). \\
$\Vsub$ as algebraic substrate ($E_8 \to H_4$ projection, icosian closure)
& Imported (Theorem)
& \cite{CascadeAlgSubstrate}: \texttt{thm:pi-H} (projection of shell onto $\Vsub$, not $E_8$ action), \texttt{thm:icosian-closure}. \\
$\Vsub$ Euclidean / conjugacy shell sizes
& Exact (imported, locally reproduced numerically)
& Exact via \cite{CascadeAlgSubstrate}~\texttt{thm:shell-class}; floating-point reproduction in Appendix~\ref{app:field-signature-audit}, entry~E. \\
$\Vsub$ nine-eigenvalue Laplacian spectrum and $94/26$ $\sigma$-fixed / $\sigma$-paired split
& Imported character-table computation, locally reproduced numerically
& Spectrum imported from \cite{CascadeAlgSubstrate} (standard $2I$ character-table); floating-point reproduction in Appendix~\ref{app:field-signature-audit}, entry~A. \\
$\Vsub$ as the unique cascade substrate
& Not claimed
& Out of scope for this paper. \\
$\|\Cphi^{-1}\| = \pphi^2$ exactly
& Analytic identity, locally reproduced
& Follows from $\lambda_{\min}(L_M) = 0$ and the $\pphi^{-2}$ shift; numerical reproduction in Appendix~\ref{app:field-signature-audit}, entry~E. \\
Observer-process tuple $\Ocal = \langle B,S,M,G,I,C,A\rangle$
& Definition~\ref{def:observer-process}
& Architectural definition; ARIA is mapped to the architecture at the process-description level, without thereby proving consciousness. \\
ARIA empirical witness ($17/18$ under standard validation; $18/18$ after documented methodology refinement; $6/6$ on a single deterministic seed; HCP $-11.58\sigma$ at ARIA $|V| = 120$ vs HCP ICA-50 $|V| = 50$)
& Imported empirical fact
& \cite{SmartARIAChess}; substrate-witness scope only. The cited source notes that the source's test P10 was validated with $15$ permutations rather than the originally preregistered $20$, with the prereg-exact rerun left open; this caveat is carried forward here. \\
$b$-anomaly closure-kernel witness
& Externally reported, commit-pinned, not locally re-executed
& \cite{SmartBAnomaly}; AIC-inconclusive ($w_{\mathrm{FREE\_C9}} = 0.652$ vs $w_{\mathrm{VFD}} = 0.348$, not a VFD preference); commit-pinned in the bibliography to \texttt{6bdfc8f8a7da9056ffeb9137d03627d06399bccf}, but external to this repository and not re-executed in this paper. \\
Standalone mass-spectrum, brain/EEG, visible-universe, Millennium-class applications (beyond the cited substrate witnesses)
& Out of scope
& Each belongs in its own paper. \\
Consciousness / mind / qualia derivation
& Not claimed
& \S\ref{ssec:claim-boundary}. \\
Replacement for quantum mechanics
& Not claimed
& \S\ref{sec:motivation}. \\
Refutation of Penrose Objective Reduction
& Not claimed
& \S\ref{sec:motivation}; OR is the structural neighbour, not the opponent. \\
\bottomrule
\end{longtable}

\subsection{What this paper consolidates}

A single mechanism name --- cascade closure --- for an operational
pattern of geometry-coupled compatibility propagation: given a
supplied closure-compatible fixed point at the top rung, the
mechanism specifies the conditions under which residual zero and
fixed-point status are propagated to the projected tower (and projection-compatibility holds by construction, since the projected tower is defined as the downward projection of the top-rung fixed point).
Selection of the top fixed point itself is not part of the present
paper. An explicit \emph{consumer API}
(\S\ref{ssec:consumer-api}): four obligations (O0)--(O3) successor
papers must discharge to invoke cascade closure as a load-bearing
premise, together with a worked status report
(\S\ref{ssec:worked-invocation}) in the programme's most
load-bearing imported tower (the finite-level P3/P4 refinement
spaces plus closure dynamics), where the source-side companion
\cite{CascadeRefinementCompat} closes (O3) and discharges the
boundary-energy / Schur substitute for (O2) in an abstract scalar
pure-midpoint refinement model, leaves (O2) as stated false, and
reduces (O1) to selection of a nontrivial element of the connected
graph Laplacian's kernel; the P3/P4-tower lift is conditional on
the open hypotheses (L1) and (L3) of
\cite[\S\texttt{sec:scope-and-gap}]{CascadeRefinementCompat}. A compressed nine-row programme-placement table
(Appendix~\ref{app:rung-placement}) showing where cascade closure
sits as the dependency gate between imported substrate rows and
downstream witness rows. A finite, audit-anchored geometric carrier
($\Vsub$ together with $\Cphi$) on which the mechanism is partly
explicit, conditional on the stated residual, projection, and
intertwining hypotheses. An
architectural identification of an existing computational system
(ARIA) as an architecture-level observer-process mapping, accompanied by the
caveated externally preregistered substrate-witness profile of
\cite{SmartARIAChess}. One short physical
witness ($b$-anomaly closure kernel) connecting the substrate operator
to an external observable.

\subsection{What is out of scope}

The following are not claimed in this paper. A unique substrate
theorem for $\Vsub$. A derivation of the visible universe from cascade
primitives. A derivation of the Standard-Model mass spectrum (the
spectral projection-class material documented in
\cite{SmartV, SmartVRev} sits in its own preprints; we cite those as
out-of-scope context only and make no statement about reciprocal
citation). A derivation of consciousness from substrate dynamics. A
resolution of any open mainstream-mathematics problem; the
Millennium-class projection scaffolding belongs in its own paper. A
generalised cosmogenesis from substrate primitives.

\subsection{Closing}

Cascade closure is one name for one pattern: in systems satisfying
the stated residual, projection, and intertwining hypotheses, if a
cascade-closure event in the sense of
Definition~\ref{def:cascade-closure} occurs, the stipulated
resolution is selection of a closure-compatible fixed point with
zero rungwise residual and projection-compatible terminal tower.
That is the conditional template defined here. The paper does no more.

\appendix

\section{Field-signature audit}\label{app:field-signature-audit}

This appendix records the load-bearing data claims of the paper. The
$\sigma$-attractor spectrum (entry A) and closure kernel (entry E)
have local numerical re-runs from this repository; the ARIA witness
profile (entry F) is an imported empirical witness from
\cite{SmartARIAChess}, not re-executed here. Re-run 2026-05-02.
The lettering A, E, F is preserved from the full
\texttt{papers/cascade-mechanism/FIELD\_SIGNATURE\_AUDIT.md} (which
also contains entries B, C, D for the mass-spectrum, octonion
derivation, and pentagonal-clock builds that are out of scope for
this paper); the gap is intentional, not a typo.

\phantomsection
\subsection*{A. \texorpdfstring{$\sigma$}{sigma}-attractor spectrum}\label{app:fsa-A}
Script: \texttt{papers/cascade-derivation/scripts/test\_sigma\_attractor\_spectrum.py}.
Graph Laplacian on $\Vsub$ has nine distinct eigenvalues: five
integer ($\{0, 9, 12, 14, 15\}$, $\sigma$-fixed under
$\sigma(\pphi) = 1 - \pphi$) and four irrational in $\Q(\sqrt{5})$
($\sigma$-paired). Counted with multiplicity: $\sigma$-fixed fraction
$94/120 = 78.3\%$; $\sigma$-paired fraction $26/120 = 21.7\%$.
Honest-negative: the $\sigma$-fixed spectral subspace has average
$2T$-mass at baseline (basis-invariant projector trace ratio; target
$\geq 1.5\times$ baseline; observed $1.0\times$). The $\sigma$-fixed
structure is purely spectral, not spatial. The
spectral-mode $\to$ Mellin-zero correspondence is conjectural
(\texttt{conj:spec\_to\_zero} in the source repository) and is not
invoked by this paper.

\phantomsection
\subsection*{E. Closure kernel}\label{app:fsa-E}
Script: \texttt{papers/aria-closure-kernel/repro/verify\_kernel.py}.
The graph $\Vsub$ is uniform-degree $12$ (recorded in the script's
\texttt{results.json}). Euclidean / conjugacy shell sizes
$[1, 12, 20, 12, 30, 12, 20, 12, 1]$ exact;
$\|\Cphi^{-1}\| \approx 2.618034$ matches $\pphi^2$ to machine precision;
per-vertex correlation uniform across all $120$ source-vertex sweeps (max $-$ min $= 2 \times
10^{-15}$). Honest-negative: UNWEIGHTED variant wins at per-vertex
correlation $0.976$ vs PHI\_GEOMETRIC $0.888$ and PHI\_ARITHMETIC
$0.884$. The two tested $\pphi$-cocycle weighted variants underperform the unweighted operator; do not claim weighted
substrate selection.

\phantomsection
\subsection*{F. ARIA observer-process witness}\label{app:fsa-F}
Source: \texttt{papers/aria-chess-paper/paper/main.tex} and the
included section files. Reviewer anchors: $17/18$ tally and
methodology refinement at
\texttt{sections/05\_results.tex}; six drug/sleep EEG signatures at
\texttt{sections/04\_consciousness\_chain.tex} and
\texttt{sections/05\_results.tex}; HCP $-11.58\sigma/+6.80\sigma$ and
the P10 $15$-vs-$20$ permutation caveat at
\texttt{sections/05\_results.tex} and
\texttt{sections/09\_limitations.tex}. Eighteen
preregistered cortical correspondences hold at $17/18$ under standard
validation and $18/18$ after the documented methodology refinement,
all at unchanged thresholds. Caveat carried over from the cited
ARIA source: the source's test P10 was validated with $15$
permutations rather than the originally preregistered $20$, with
the prereg-exact rerun left open in the source.
Six drug/sleep EEG signatures pass on a single deterministic
recurrent-layer trajectory at seed $42$. HCP brain functional
connectivity degree-homogeneity (degree standard-deviation) places
ARIA at $-11.58\sigma$ vs HCP baseline (ARIA $|V| = 120$ vs HCP ICA-50
$|V| = 50$); the clustering-coefficient comparison is a separate
$+6.80\sigma$ result and should not be conflated with the homogeneity
figure.

\subsection*{Re-run command}

The canonical reproduction artefact for entries~A and~E is the
self-contained bundle at
\texttt{papers/cascade-mechanism/repro/}, which calls both scripts,
captures their output, and diffs against frozen reference outputs
shipped in \texttt{papers/cascade-mechanism/repro/expected\_outputs/}.
Reviewers should run, from a writable working tree:

\begin{verbatim}
cd <repo-root>/papers/cascade-mechanism/repro
bash run_audit.sh
\end{verbatim}

A clean run prints \texttt{AUDIT PASSED}. The bundle's
\texttt{README.md} documents which Appendix entry each script
anchors and what is intentionally not reproduced inside the bundle
(ARIA preregistered profile and $b$-anomaly fit; both external,
both pinned in the bibliography). Underlying scripts (unchanged):
\texttt{papers/cascade-derivation/scripts/test\_sigma\_attractor\_spectrum.py}
and \texttt{papers/aria-closure-kernel/repro/verify\_kernel.py};
\texttt{verify\_kernel.py} writes
\texttt{papers/aria-closure-kernel/repro/results.json}.

ARIA witness (entry F) is documented in
\texttt{papers/aria-chess-paper/paper/main.tex}; the upstream re-run
lives in the \texttt{ariaChessRepo}~\cite{ariaChessRepo}.

\section{ARIA runtime module mapping}\label{app:aria-runtime-map}

The observer-process tuple of Definition~\ref{def:observer-process} is
mapped to ARIA source tiers across two layers. Tier~1 is the in-repo
reproducibility bundle at
\texttt{release-bundle/closure-kernel-papers/papers/aria-chess/repro/kernel/},
which contains six substantive \texttt{.py} modules (excluding \texttt{\_\_init\_\_.py})
(\texttt{vfd\_closure\_kernel.py}, \texttt{self\_model\_stream.py},
\texttt{lyapunov\_selector.py}, \texttt{sigma\_orbit\_basis.py},
\texttt{v66\_e8\_coxeter.py}, and
\texttt{consciousness\_binding.py}). Five of these are mapped to
$\Ocal$ components in the table below; the sixth,
\texttt{consciousness\_binding.py}, is in the bundle but is
deliberately left unmapped here for the same reason that consciousness
or qualia derivation is out-of-scope for this paper. Note that
\texttt{consciousness\_binding.py} remains a Tier~1 \emph{runtime}
dependency of \texttt{self\_model\_stream.py} (i.e.\ the bundle does
not function with that file removed); the choice not to assign it an
$\Ocal$-component label is a scope decision of the present paper, not
a claim about the bundle's runtime structure. The bundle is
sufficient for the kernel/demo-level numerical reproductions in
\cite{AriaClosureKernel}; it is \emph{not} sufficient to re-execute the
full preregistered validation suite of \cite{SmartARIAChess} (which
requires the upstream brain-validation harnesses, demo scripts, and
the v4 substrate parameters frozen in the Tier~2 snapshot). Tier~2 is the upstream production runtime
\texttt{ariaChessRepo}~\cite{ariaChessRepo}, which contains the full
production observer-process architecture represented at the module-description level across boundary
sampling, state integration, memory, governance, closure, action, and
replay. The mapping below was constructed 2026-05-02. Tier~1 files
exist in the in-repo bundle and are locally verifiable. Tier~2 entries
are external upstream witnesses, pinned at the directory level to the
frozen working-tree snapshot \texttt{v4\_locked\_2026-04-29/} of
\cite{ariaChessRepo}. The upstream working tree is not under
\texttt{git} version control, so a reviewer-accessible cryptographic
pin (e.g.\ a git commit hash) is not available; in lieu of that, the
snapshot is identified through the directory-level artefacts
documented in the bibliography entry (snapshot date 2026-04-29,
session UUID \texttt{97c7fa0b-8d38-465b-a769-26c9983b446d}, substrate
version v4 with \texttt{k\_threshold=12}, deterministic seeds 42 and
32000--32019, and an author-side off-repo tarball at
\texttt{\textasciitilde{}/.aria/snapshots/aria-chess-v4-2026-04-29.tar.gz}
that is not by itself reviewer-accessible). Tier~2 module-level
claims are external upstream evidence at the file/path level; the
present paper does not re-execute the upstream production runtime.

\renewcommand{\arraystretch}{1.30}
\begin{longtable}{>{\raggedright\arraybackslash}p{0.18\textwidth}>{\raggedright\arraybackslash}p{0.29\textwidth}>{\raggedright\arraybackslash}p{0.43\textwidth}}
\caption{Two-tier runtime mapping of $\Ocal = \langle B, S, M, G, I, C, A\rangle$
to ARIA modules.}\label{tab:runtime-map}\\
\toprule
\textbf{$\Ocal$ component} & \textbf{Tier 1 (in-repo bundle)} & \textbf{Tier 2 (upstream \texttt{ariaChessRepo})} \\
\midrule
\endfirsthead
\toprule
\textbf{$\Ocal$ component} & \textbf{Tier 1 (in-repo bundle)} & \textbf{Tier 2 (upstream \texttt{ariaChessRepo})} \\
\midrule
\endhead
$B$ boundary sampler
& \texttt{vfd\_closure\_kernel.py}
& \texttt{raw\_sensory.py}; \texttt{action\_encoder.py}; \texttt{chess\_runtime\_adapter.py}\\
$S$ state integrator
& \texttt{self\_model\_stream.py}
& \texttt{coherence\_runtime.py}; \texttt{coherence\_selfmodel.py}; \texttt{self\_model\_stream.py}\\
$M$ memory substrate
& \texttt{self\_model\_stream.py} state EMAs
& \texttt{coherence\_memory.py}; \texttt{polytope\_persistence.py}; \texttt{coherence\_longhorizon.py}\\
$G$ goal/context
& runtime knobs in \texttt{v66\_e8\_coxeter.py}
& \texttt{coherence\_planner.py}; \texttt{coherence\_selfmodel.py}\\
$I$ invariant constraints
& --- (not in bundle)
& \texttt{coherence\_constitution.py}; \texttt{coherence\_audit.py}; \texttt{coherence\_review.py}\\
$C$ closure/crystallisation
& \texttt{vfd\_closure\_kernel.py}; \texttt{lyapunov\_selector.py}; \texttt{sigma\_orbit\_basis.py}; \texttt{v66\_e8\_coxeter.py}
& \texttt{coherence\_crystallisation.py}; \texttt{vfd\_closure\_kernel.py}; \texttt{lyapunov\_selector.py}; \texttt{dual\_closure\_operator.py}\\
$A$ action interface
& --- (not in bundle)
& \texttt{coherence\_execution.py}; \texttt{coherence\_delivery\_calibration.py}\\
\bottomrule
\end{longtable}

The mapping is given by audited file location, not by reverse
inference from runtime behaviour. Components $I$ and $A$ are not
present in Tier~1 and are present in Tier~2 only; this is a fact
about the in-repo bundle, not a limitation of the architecture.

\section{Programme placement of cascade closure}\label{app:rung-placement}

This appendix records, in one table, where the cascade-closure
mechanism of \S\ref{sec:mechanism} sits relative to the other named
geometric layers of the present programme. No row introduces a new
mathematical claim; every row imports its substantive content from a
named source paper, with the source's own labels preserved. The table
is intended as a navigation aid for successor papers invoking
cascade closure under the consumer API of
\S\ref{ssec:consumer-api}, not as a programme survey.

Status flags: \textsf{[Th-imp]} imported theorem; \textsf{[Pr-imp]}
imported proposition; \textsf{[Hyp-imp]} imported under a named
conditional hypothesis (label tracked inside the source); \textsf{[Mod]}
model choice within the programme, not derived elsewhere as a
uniqueness theorem; \textsf{[Int]} interpretive use of structural
material, not load-bearing for any formal claim in the present
paper. A row's status is determined by the cited source, not by the
present paper.

\renewcommand{\arraystretch}{1.20}
\begin{longtable}{>{\raggedright\arraybackslash}p{0.20\textwidth}>{\raggedright\arraybackslash}p{0.32\textwidth}>{\raggedright\arraybackslash}p{0.10\textwidth}>{\raggedright\arraybackslash}p{0.30\textwidth}}
\caption{Programme placement of cascade closure. Compressed
nine-row version of the v1 maximalist rung table, retained here as
the consumer-API navigation aid.}\label{tab:rung-placement}\\
\toprule
\textbf{Layer} & \textbf{Geometry / object} & \textbf{Status} & \textbf{Primary source} \\
\midrule
\endfirsthead
\toprule
\textbf{Layer} & \textbf{Geometry / object} & \textbf{Status} & \textbf{Primary source} \\
\midrule
\endhead
Arithmetic coefficients
& $\Q(\pphi)$, $\Z[\pphi]$; coefficientwise Galois involution $\sigma : \pphi \leftrightarrow 1 - \pphi = -1/\pphi$
& \textsf{[Th-imp]}
& \cite{CascadeSigmaRat} (definition and fixed-part theorems for $\sigma$) \\
Algebraic substrate
& $E_8$; icosian $H_4$; binary icosahedral $2I$; 600-cell $\Vsub$; $D_4$
& \textsf{[Th-imp]} / \textsf{[Pr-imp]}
& \cite{CascadeAlgSubstrate}: \texttt{thm:pi-H} (projection, not action), \texttt{thm:icosian-closure}, \texttt{thm:shell-class} \\
Refinement state spaces and bonding maps
& Finite $X_n$; bonding maps $p_{n+1, n}$
& \textsf{[Th-imp]}
& \cite{CascadeRefinementSpaces}: \texttt{thm:bonding}; specific $p^0, p^1$ identities under \texttt{thm:commute} (not generic) \\
Closure dynamics (gradient flow infrastructure)
& Quadratic finite-level $F_n = \alpha R_n + \beta E_n - \gamma Q_n$; gradient operator $L_n$; flow $e^{-tL_n}$
& \textsf{[Th-imp]}
& \cite{CascadeClosureDynamics}: \texttt{thm:gradient-operator}, \texttt{prop:energy-dissipation}, \texttt{thm:coercive-contraction}, \texttt{prop:mass-refinement}, \texttt{thm:flow-exists}, \texttt{thm:flow-intertwining} (mass-only $\alpha = \beta = 0$) \\
\textbf{Cascade-closure mechanism (this paper)}
& Definition~\ref{def:cascade-closure}; Proposition~\ref{prop:resid-comp}; Corollary~\ref{cor:conv-prop}; consumer API \S\ref{ssec:consumer-api}
& --- \textbf{(gate)}
& Imported substrate from rows above; gates rows below by supplying the named selection-mechanism prerequisite. Not load-bearing for itself \\
Passive response kernel
& $\Cphi = L_M + \pphi^{-2} I$ on $\Vsub$; $\|\Cphi^{-1}\| = \pphi^2$ exactly; $b$-anomaly external substrate witness
& \textsf{[Th-imp]} / \textsf{[Mod]}
& \cite{AriaClosureKernel}; \cite{SmartBAnomaly} (commit-pinned, externally reported, not re-executed here) \\
Spectral mass projection class
& $\Vsub$ Laplacian eigenvalues / shell decomposition; mass-coefficient assignment map (out-of-scope context for this paper)
& \textsf{[Mod]}
& \cite{SmartV, SmartVRev} \\
Active / cognitive substrate
& $H_4$ / 600-cell active substrate mapping; preregistered EEG / chess / HCP substrate-witness profile (architecture-level mapping, see \S\ref{sec:aria})
& \textsf{[Mod]}
& \cite{SmartARIAChess, ariaChessRepo} \\
Observer-process architecture
& $\Ocal = \langle B, S, M, G, I, C, A\rangle$ (Definition~\ref{def:observer-process})
& \textsf{[Mod]}
& Defined here (Definition~\ref{def:observer-process}); \cite{ObserverInstance} is related upstream context (a five-tuple slot definition, not the seven-tuple architecture used here), and \S\ref{sec:aria} gives the runtime mapping. \\
\bottomrule
\end{longtable}

The single load-bearing observation in this table is the
\textbf{(gate)} row: cascade closure imports its analytic
infrastructure entirely from the rows above it (arithmetic,
algebraic substrate, refinement spaces, closure dynamics) and
gates the rows below it (passive kernel, mass projection, active
substrate, observer-process architecture) by supplying the named
selection-mechanism prerequisite they need to invoke. The rows
below cite the gate row through the consumer API of
\S\ref{ssec:consumer-api}; the rows above are unaffected by the
present paper.

\bibliographystyle{plain}
\bibliography{references}

\end{document}
